\documentclass[../../main.tex]{subfiles}
\begin{document}

\begin{center}
  {\fontsize{48pt}{58pt}\selectfont \textbf{2002}\normalsize\textbf{年度}}
\end{center}

\begin{table}[htb]
  \centering
  \caption{出題テーマと難易度}
  \begin{tabular}{|c|p{6.2cm}|c|c|c|}
    \hline
    \textbf{問題番号} & \textbf{テーマ} & \textbf{思考量} & \textbf{計算量} & \textbf{総合} \\
    \hline
    第1問 & 二つの放物線が交わる条件 & 易 & 普 & 易 \\
    \hline
    第2問 & 整式の除法と数列の漸化式 & 普 & 易 & 易 \\
    \hline
    第3問 & 球面から定まる空間領域と体積 & 普 & 普 & 普 \\
    \hline
    第4問 & 曲線の接線と直線の直交条件 & 易 & 普 & 普 \\
    \hline
    第5問 & 三角錐の体積の和と区分求積 & 易 & 普 & 易 \\
    \hline
    第6問 & 数列のシャッフルと合同式 & 難 & 普 & 普 \\
    \hline
  \end{tabular}
\end{table}

\begin{table}[htb]
  \centering
  \caption{解答}
  \renewcommand{\arraystretch}{1.4}
  \begin{tabular}{|c|c|p{9cm}|}
    \hline
    \textbf{問題番号} & \textbf{小問} & \textbf{答え} \\
    \hline
    第1問 & -- & $\dfrac{4\pi}{3}+2n\pi<\theta<
    \dfrac{5\pi}{3}+2n\pi$（$n\in\mathbb{Z}$） \\
    \hline
    第2問 & (1) & $a_{n+1}=a_n+b_n$，$b_{n+1}=a_n$（証明問題） \\
    \cline{2-3}
    & (2) & $a_n,b_n$は正の整数で互いに素（証明問題） \\
    \hline
    第3問 & -- & 領域$V$は解答の図を参照．
    体積は$\dfrac{19\pi}{6}<10$ \\
    \hline
    第4問 & -- & $a>\dfrac{1}{2}$ \\
    \hline
    第5問 & -- & $\displaystyle\lim_{n\to\infty}
    \sum_{k=0}^{n-1}V_k=\dfrac{\sqrt{2}\pi}{48}$ \\
    \hline
    第6問 & (1) & $\{8,7,6,5,4,3,2,1\}$ \\
    \cline{2-3}
    & (2) & $f(k)-2k$は$2N+1$で割り切れる（証明問題） \\
    \cline{2-3}
    & (3) & $N=2^{n-1}$なら$2n$回のシャッフルで元に戻る
    （証明問題） \\
    \hline
  \end{tabular}
\end{table}

\paragraph{試験形態}
6問．試験時間は要確認．

\end{document}
